\bta{2001}


\begin{enumerate}
	%\renewcommand{\labelenumi}{\arabic{enumi}.}
	% A(\Alph) a(\alph) I(\Roman) i(\roman) 1(\arabic)
	%设定全局标号series=example	%引用全局变量resume=example
	%[topsep=-0.3em,parsep=-0.3em,itemsep=-0.3em,partopsep=-0.3em]
	%可使用leftmargin调整列表环境左边的空白长度 [leftmargin=0em]
	\item
填空题 (每小题 3 分，共 15 分)
\begin{enumerate}
	%\renewcommand{\labelenumi}{\arabic{enumi}.}
	% A(\Alph) a(\alph) I(\Roman) i(\roman) 1(\arabic)
	%设定全局标号series=example	%引用全局变量resume=example
	%[topsep=-0.3em,parsep=-0.3em,itemsep=-0.3em,partopsep=-0.3em]
	%可使用leftmargin调整列表环境左边的空白长度 [leftmargin=0em]
	\item
$F(x)=\frac{x^{2}}{x-a} \int_{a}^{x} f(t) d t$, 其中 $f(x)$ 连续, 则有 $\lim\limits _{x \rightarrow a} F(x)$ \tk{$a^{2}f(a)$} 。
\item 
通过点 $(3,1,-1)$ 及直线 $\frac{x-4}{5}=\frac{y+3}{2}=\frac{z}{1}$ 的平面方程为: \tk{$-3x+2y+11z+18=0$} 。

\item 
$\lim\limits _{x \rightarrow 0} \frac{\sqrt{1+x}+\sqrt{1-x}-2}{x^{2}}=$ \tk{$-\frac14$} 。

\item 
方程 $y^{\prime \prime}+y^{\prime}+y=0$ 的通解为
\tk{$e^{-x/2}\left(C_{1}\cos\frac{\sqrt3x}{2}+C_{2}\sin\frac{\sqrt3x}{2}\right)$} 。


\item 
设方程 $f\left(\frac{z}{x}, \frac{y}{z}\right)=0$ 确立了隐函数 $z=z(x, y)$, 则 $\frac{\partial z}{\partial x}=$
\tk{$\displaystyle\frac{z^{3}f_{1}}{x\left(z^{2}f_{1}-xyf_{2}\right)}$}，
其中 $f_{1},f_{2}$ 均在 $\left(\frac zx,\frac yz\right)$ 处取值。
	
	
	
	
\end{enumerate}

\banswer{
	\begin{enumerate}
		%\renewcommand{\labelenumi}{\arabic{enumi}.}
		\item
		$a^{2}f(a)$。由变上限积分的平均值极限（积分中值定理）
		\[
		\lim_{x\to a}\frac{1}{x-a}\int_a^x f(t)\,dt=f(a),
		\]
		又 $x^2\to a^2$，故
		\[
		\lim_{x\to a}F(x)
		=\lim_{x\to a}x^2\cdot\frac{\int_a^x f(t)\,dt}{x-a}
		=a^2f(a).
		\]

		\item
		$-3x+2y+11z+18=0$。所给直线经过 $Q=(4,-3,0)$，方向向量为
		$\bm d=(5,2,1)$。设指定点为 $P=(3,1,-1)$，则
		\[
		\overrightarrow{QP}=(-1,4,-1),\qquad
		\bm d\times\overrightarrow{QP}=(-6,4,22)
		=2(-3,2,11).
		\]
		因此可取平面的法向量 $\bm n=(-3,2,11)$。代入点 $P$，
		平面方程为
		\[
		-3(x-3)+2(y-1)+11(z+1)=0,
		\]
		即 $-3x+2y+11z+18=0$。
		（校验：$Q=(4,-3,0)$ 与直线上另一点 $(9,-1,1)$ 都满足此式。）

		\item
		$-\frac14$。记 $A=\sqrt{1+x}+\sqrt{1-x}$。连续两次有理化可得
		\begin{align*}
		A-2
		&=\frac{A^2-4}{A+2}
		=\frac{2\bigl(\sqrt{1-x^2}-1\bigr)}{A+2}\\
		&=-\frac{2x^2}
		{(A+2)\bigl(\sqrt{1-x^2}+1\bigr)}.
		\end{align*}
		所以
		\[
		\lim_{x\to0}\frac{A-2}{x^2}
		=-\frac{2}{(2+2)(1+1)}
		=-\frac14 .
		\]

		\item
		$y=e^{-x/2}\left(C_1\cos\frac{\sqrt3x}{2}+C_2\sin\frac{\sqrt3x}{2}\right)$。
		特征方程为
		\[
		r^2+r+1=0,\qquad
		r_{1,2}=-\frac12\pm\frac{\sqrt3}{2}i ,
		\]
		是一对共轭复根 $\alpha\pm\beta i$（$\alpha=-\frac12$，$\beta=\frac{\sqrt3}{2}$），
		故实通解如上。

		\item
		$\dfrac{z^3f_1}{x\left(z^2f_1-xyf_2\right)}$。令
		\[
		u=\frac zx,\qquad v=\frac yz,
		\]
		并记 $f_1=f_u(u,v)$、$f_2=f_v(u,v)$。求
		$\partial z/\partial x$ 时把 $y$ 看作常量。对
		$f(u,v)=0$ 关于 $x$ 求偏导：
		\[
		f_1\left(\frac{z_x}{x}-\frac{z}{x^2}\right)
		+f_2\left(-\frac{y z_x}{z^2}\right)=0.
		\]
		整理为
		\[
		\left(\frac{f_1}{x}-\frac{yf_2}{z^2}\right)z_x
		=\frac{zf_1}{x^2},
		\]
		在分母不为零的条件下，
		\[
		\frac{\partial z}{\partial x}
		=\frac{z^3f_1}{x\left(z^2f_1-xyf_2\right)} .
		\]
	\end{enumerate}
}




\item 
(本题满分 10 分)
设 $y=y(x), z=z(x)$ 由方程 $z=x f(x+y)$ 和 $F(x, y, z)=0$ 所确定的函数, 其中 $f$ 和 $F$ 分别具有一阶连续导数和一阶连续偏导数, 求 $\frac{d z}{d x}$。



\banswer{
	记 $f=f(x+y)$，$f'=f'(x+y)$，并将 $F_x,F_y,F_z$ 均理解为在
	$(x,y,z)$ 处的偏导数。对两个方程求导得
	\[
	z'=f+xf'(1+y'),\qquad F_x+F_yy'+F_zz'=0.
	\]
	把第一式写成
	\[
	z'-xf'y'=f+xf',
	\]
	与第二式组成关于 $y'$、$z'$ 的线性方程组
	\[
	\begin{cases}
		-xf'y'+z'=f+xf',\\
		F_yy'+F_zz'=-F_x.
	\end{cases}
	\]
	由第一式 $z'=f+xf'+xf'y'$，代入第二式，得到
	\[
	(F_y+xf'F_z)y'=-F_x-(f+xf')F_z.
	\]
	再代回 $z'$ 并合并同类项，在
	$F_y+xf'F_z\ne0$ 时有
	\[
	\boxed{\displaystyle
	\frac{dz}{dx}
	=\frac{(f+xf')F_y-xf'F_x}{F_y+xf'F_z}} .
	\]
}


\item 
(本题满分 12 分)
\begin{enumerate}
	%\renewcommand{\labelenumi}{\arabic{enumi}.}
	% A(\Alph) a(\alph) I(\Roman) i(\roman) 1(\arabic)
	%设定全局标号series=example	%引用全局变量resume=example
	%[topsep=-0.3em,parsep=-0.3em,itemsep=-0.3em,partopsep=-0.3em]
	%可使用leftmargin调整列表环境左边的空白长度 [leftmargin=0em]
	\item
(5分)
计算
	\[
	\int_{0}^{\pi} \frac{\sin \theta}{\sqrt{1-2 a \cos \theta+a^{2}}}\,d\theta,\quad(a>1)
	\]
	
	
	\item
	(7分)
	计算曲面积分
	\[
	I=\iint_{S}\left(x^{2}+y^{2}\right) d S
	\]
	$S$ : 曲面 $z=\sqrt{x^{2}+y^{2}}$ 及 $z=1$ 所围成立体的全表面。
	
	
	
	
\end{enumerate}

\banswer{
	\begin{enumerate}
		\item
		令 $u=1-2a\cos\theta+a^2$，则
		$du=2a\sin\theta\,d\theta$。当 $\theta=0,\pi$ 时，
		$u$ 分别为 $(a-1)^2,(a+1)^2$。因为 $a>1$，
		$\sqrt{(a\pm1)^2}=a\pm1$，所以
		\[
		\begin{aligned}
		\int_0^\pi\frac{\sin\theta}
		{\sqrt{1-2a\cos\theta+a^2}}\,d\theta
		&=\frac1{2a}\int_{(a-1)^2}^{(a+1)^2}u^{-1/2}\,du\\
		&=\frac1a\left[\sqrt{u}\right]_{(a-1)^2}^{(a+1)^2}\\
		&=\frac1a\bigl[(a+1)-(a-1)\bigr]
		=\boxed{\frac2a}.
		\end{aligned}
		\]
		\item
		全表面由圆锥侧面 $S_1$ 和顶面圆盘 $S_2$ 组成。
		圆锥面可参数化为
		\[
		\bm r(r,\theta)=(r\cos\theta,r\sin\theta,r),
		\quad 0\leqslant r\leqslant1,\quad
		0\leqslant\theta\leqslant2\pi.
		\]
		由
		\[
		\left|\bm r_r\times\bm r_\theta\right|=\sqrt2\,r
		\]
		知 $dS_1=\sqrt2\,r\,dr\,d\theta$。在顶面
		$z=1$ 上，极坐标给出 $dS_2=r\,dr\,d\theta$。
		两部分均有 $x^2+y^2=r^2$，故
		\[
		\begin{aligned}
		I=2\pi\sqrt2\int_0^1r^3\,dr
		+2\pi\int_0^1r^3\,dr
		&=\frac{\pi\sqrt2}{2}+\frac{\pi}{2}\\
		&=\boxed{\frac{\pi}{2}(1+\sqrt2)}.
		\end{aligned}
		\]
	\end{enumerate}
}



\item 
(本题满分 8 分)
将函数 $f(x)=\arctan \frac{1+x}{1-x}$ 展开为 $x$ 的幂级数。

\banswer{
	利用正切的加法公式，
	\[
	\tan\left(\frac{\pi}{4}+\arctan x\right)
	=\frac{1+x}{1-x}.
	\]
	当 $|x|<1$ 时，两边对应的角均落在反正切函数的主值范围内，
	所以
	\[
	\arctan\frac{1+x}{1-x}
	=\frac{\pi}{4}+\arctan x.
	\]
	又由几何级数
	\[
	\frac1{1+t^2}=\sum_{n=0}^{\infty}(-1)^nt^{2n},
	\qquad |t|<1,
	\]
	在 $0$ 到 $x$ 上逐项积分，得到
	\[
	\boxed{\displaystyle
	\arctan\frac{1+x}{1-x}
	=\frac{\pi}{4}
	+\sum_{n=0}^{\infty}\frac{(-1)^n}{2n+1}x^{2n+1}},
	\qquad |x|<1.
	\]
	该幂级数的收敛半径为 $R=1$。在 $x=-1$ 时级数收敛到
	$0$，且原函数也有定义；在 $x=1$ 时级数本身收敛到
	$\pi/2$，但原式分母为零，故原函数无定义。
}



\item 
(本题满分 20 分)
今有常微分方程组
\[
\left\{\begin{array}{l}
	\frac{d x}{d t}=3 x+2 y-z \\
	\frac{d y}{d t}=-2 x-2 y+2 z \\
	\frac{d z}{d t}=3 x+6 z-z
\end{array}\right.
\]	
\begin{enumerate}
	%\renewcommand{\labelenumi}{\arabic{enumi}.}
	% A(\Alph) a(\alph) I(\Roman) i(\roman) 1(\arabic)
	%设定全局标号series=example	%引用全局变量resume=example
	%[topsep=-0.3em,parsep=-0.3em,itemsep=-0.3em,partopsep=-0.3em]
	%可使用leftmargin调整列表环境左边的空白长度 [leftmargin=0em]
	\item
改写方程组为 $\frac{d}{d t} X(t)=A X(t)$ 形式，其中 $X(t)=(x(t), y(t), z(t))^{\prime}$ 为列向量, $A$为 3 阶方阵, 记号 “$  {}^{\prime}  $” 表示转置。(2分)
	
	
	\item
	求矩阵 $A$ 的特征值及对应的特征向量。(10分)
	
	
	\item
	给出非退化矩阵 $T$ 和对角阵 $D$，使得 $T^{-1} A T=D$。(3分)
	
	
	\item
	给出此微分方程的通解。(5分)
	
	
	
	
\end{enumerate}


\banswer{
	\begin{enumerate}
		\item
		按题面第三式
		$\frac{dz}{dt}=3x+6z-z=3x+5z$，所以
		\[
		X'(t)=AX(t),\qquad
		A=\begin{pmatrix}
		3&2&-1\\
		-2&-2&2\\
		3&0&5
		\end{pmatrix}.
		\]
		\item
		特征多项式为
		\[
		\begin{aligned}
		|\lambda I-A|
		&=\begin{vmatrix}
		\lambda-3&-2&1\\
		2&\lambda+2&-2\\
		-3&0&\lambda-5
		\end{vmatrix}\\
		&=(\lambda-3)(\lambda+2)(\lambda-5)
		+7\lambda-26\\
		&=\lambda^3-6\lambda^2+6\lambda+4\\
		&=(\lambda-2)(\lambda^2-4\lambda-2).
		\end{aligned}
		\]
		故三个特征值为
		\[
		\lambda_1=2,\qquad
		\lambda_2=2+\sqrt6,\qquad
		\lambda_3=2-\sqrt6.
		\]
		当 $\lambda=2$ 时，$(A-2I)\bm v=0$ 给出
		$z=-x,\ y=-x$，可取 $\bm v_1=(1,-1,-1)^{\prime}$。
		对于满足 $\lambda^2-4\lambda-2=0$ 的另两个特征值，
		方程组给出
		\[
		y=-2x,\qquad z=-(\lambda+1)x.
		\]
		令 $x=1$，便得到
		相应的一组特征值与特征向量为
		\[
		\begin{array}{c|c}
		\lambda&\text{特征向量}\\ \hline
		2&(1,-1,-1)^{\prime}\\
		2+\sqrt6&(1,-2,-3-\sqrt6)^{\prime}\\
		2-\sqrt6&(1,-2,-3+\sqrt6)^{\prime}
		\end{array}.
		\]
		\item
		把上述三个特征向量依次作为列向量，可取
		\[
		T=\begin{pmatrix}
		1&1&1\\
		-1&-2&-2\\
		-1&-3-\sqrt6&-3+\sqrt6
		\end{pmatrix},\qquad
		D=\operatorname{diag}(2,2+\sqrt6,2-\sqrt6).
		\]
		由于
		\[
		\det T=-2\sqrt6\ne0,
		\]
		矩阵 $T$ 非退化；又 $AT=TD$，所以
		$T^{-1}AT=D$。
		\item
		令 $X=TY$，则 $Y'=DY$，即三个分量分别满足
		\[
		y_1'=2y_1,\qquad
		y_2'=(2+\sqrt6)y_2,\qquad
		y_3'=(2-\sqrt6)y_3.
		\]
		因此 $Y=(C_1e^{2t},C_2e^{(2+\sqrt6)t},
		C_3e^{(2-\sqrt6)t})^{\prime}$，从而原方程组的通解为
		\[
		\begin{aligned}
		X(t)={}&C_1e^{2t}(1,-1,-1)^{\prime}\\
		&+C_2e^{(2+\sqrt6)t}(1,-2,-3-\sqrt6)^{\prime}\\
		&+C_3e^{(2-\sqrt6)t}(1,-2,-3+\sqrt6)^{\prime}.
		\end{aligned}
		\]
	\end{enumerate}
}



\item 
(本题满分 5 分)
设 $A$ 为 $m \times n$ 阶矩阵, $A$ 的秩为 $r>0$ 。证明存在秩为 $r$ 的 $m \times r$ 阶矩阵 $B$ 和秩 为 $r$ 的 $r \times n$ 阶矩阵 $C$, 使得 $A=B C$。

\banswer{
	把 $A$ 的列记为
	\[
	A=(\bm a_1,\bm a_2,\ldots,\bm a_n).
	\]
	由于 $\operatorname{rank}A=r$，可从这些列中选出 $r$ 个
	线性无关的列 $\bm b_1,\ldots,\bm b_r$，使它们构成 $A$
	的列空间的一组基。令
	\[
	B=(\bm b_1,\ldots,\bm b_r),
	\]
	则 $B$ 是 $m\times r$ 矩阵且 $\operatorname{rank}B=r$。

	对每个 $j=1,\ldots,n$，存在唯一的系数
	$c_{1j},\ldots,c_{rj}$，使
	\[
	\bm a_j=c_{1j}\bm b_1+\cdots+c_{rj}\bm b_r
	=B\begin{pmatrix}c_{1j}\\ \vdots\\ c_{rj}\end{pmatrix}.
	\]
	把这些坐标列向量依次排列，定义
	\[
	C=(c_{ij})_{r\times n}.
	\]
	逐列比较立即得到 $A=BC$。最后，
	\[
	r=\operatorname{rank}A=\operatorname{rank}(BC)
	\leqslant\operatorname{rank}C\leqslant r,
	\]
	所以 $\operatorname{rank}C=r$。这就构造出了题目要求的
	$B$ 和 $C$。
}



\item 
(本题满分 5 分)
$f(z)=u(x, y)+i v(x, y)$ 是区域 $D$ 上的非常值解析函数, 试说明 $\overline{f(z)}=u-i v$ 和 $g(z)=v+i v$ 在 $D$ 上的解析性。

\banswer{
	二者都不可能在整个区域 $D$ 上解析。

	因为 $f=u+iv$ 解析，其 Cauchy--Riemann 方程为
	\[
	u_x=v_y,\qquad u_y=-v_x. \tag{1}
	\]
	若 $\overline{f}=u-iv$ 也在 $D$ 上解析，则它还应满足
	\[
	u_x=-v_y,\qquad u_y=v_x. \tag{2}
	\]
	比较 (1)、(2)，得到
	\[
	u_x=u_y=v_x=v_y=0.
	\]
	因此 $u,v$ 在区域 $D$ 上均为常数，从而 $f$ 为常值，
	与题设矛盾。所以 $\overline f$ 在 $D$ 上不解析。

	对于 $g=v+iv$，其实部和虚部都等于 $v$。若它在 $D$
	上解析，则其 Cauchy--Riemann 方程为
	\[
	v_x=v_y,\qquad v_y=-v_x.
	\]
	两式相加、相减可得 $v_x=v_y=0$。再由 (1) 得
	$u_x=u_y=0$，仍推出 $f$ 为常值，产生矛盾。因此 $g$
	也不在整个 $D$ 上解析。

	这里的结论是“不能在区域 $D$ 上处处解析”；题设本身并不排除
	它们在某些孤立点恰好满足复可微条件。
}



\item 
(本题满分 5 分)
计算积分
\[
\int_{|z|=3}\left[\bar{z}+(z-1)^{5} \cos \frac{1}{(z-1)^{3}}\right] d z
\]
其中 $\bar{z}$ 为复数 $z$ 的共轭复数。

\banswer{
	取圆周的正向。在 $|z|=3$ 上有 $\bar z=9/z$，并且
	\[
	\begin{aligned}
	(z-1)^5\cos\frac1{(z-1)^3}
	&=(z-1)^5
	\sum_{k=1}^{\infty}
	\frac{(-1)^k}{(2k)!}(z-1)^{5-6k}\\
	&=(z-1)^5-\frac1{2(z-1)}
	+\frac1{24(z-1)^7}-\cdots.
	\end{aligned}
	\]
	因此其在 $z=1$ 处的留数为 $-1/2$。另一方面，
	\[
	\int_{|z|=3}\bar z\,dz
	=\int_{|z|=3}\frac9z\,dz=18\pi i.
	\]
	由留数定理，
	\[
	\begin{aligned}
	\int_{|z|=3}\left[\bar z+(z-1)^5
	\cos\frac1{(z-1)^3}\right]dz
	&=2\pi i\left(9-\frac12\right)\\
	&=\boxed{17\pi i}.
	\end{aligned}
	\]
}


\item 
(本题满分 12 分)
用分离变量法解定解问题。
\[
\left\{\begin{array}{ll}
		\frac{\partial^{2} u}{\partial x^{2}}+\frac{\partial^{2} u}{\partial y^{2}}=0 & (x, y) \in D, D:-\infty<x<0,0<y<\pi \\
		\left.u\right|_{x=0}=\frac{y}{\pi},\left.u\right|_{x=-\infty} \text { 有界 } & (0<y<\pi) \\
		\left.u\right|_{y=0}=0,\left.u\right|_{y=\pi}=0 & (-\infty<x<0)
\end{array}\right.
\]


\banswer{
	令 $u(x,y)=X(x)Y(y)$。代入 $u_{xx}+u_{yy}=0$，得
	\[
	\frac{X''}{X}=-\frac{Y''}{Y}=\lambda.
	\]
	于是
	\[
	X''-\lambda X=0,\qquad Y''+\lambda Y=0.
	\]
	由 $Y(0)=Y(\pi)=0$，非零解只在
	$\lambda=n^2$（$n=1,2,\ldots$）时存在，此时
	\[
	Y_n(y)=\sin(ny).
	\]
	相应的 $x$ 方程为 $X_n''-n^2X_n=0$，其通解是
	\[
	X_n(x)=A_ne^{nx}+B_ne^{-nx}.
	\]
	当 $x\to-\infty$ 时，$e^{-nx}$ 无界，所以必须令
	$B_n=0$。因此满足三条齐次条件及有界条件的级数形式为
	\[
	u(x,y)=\sum_{n=1}^{\infty}b_ne^{nx}\sin(ny).
	\]
	再利用 $x=0$ 处的边界条件，将 $y/\pi$ 展开为正弦级数：
	\[
	\frac{y}{\pi}
	=\sum_{n=1}^{\infty}b_n\sin(ny),\qquad
	\]
	其中分部积分给出
	\[
	\begin{aligned}
	b_n
	&=\frac{2}{\pi}\int_0^\pi\frac{y}{\pi}\sin(ny)\,dy\\
	&=\frac{2}{\pi^2}
	\left[-\frac{y\cos(ny)}n+\frac{\sin(ny)}{n^2}\right]_0^\pi\\
	&=\frac{2(-1)^{n+1}}{\pi n}.
	\end{aligned}
	\]
	代回即得
	\[
	\boxed{\displaystyle
	u(x,y)=\frac2\pi\sum_{n=1}^{\infty}
	\frac{(-1)^{n+1}}{n}e^{nx}\sin(ny)}.
	\]
}


\item 
(本题满分 8 分)
\begin{enumerate}
	%\renewcommand{\labelenumi}{\arabic{enumi}.}
	% A(\Alph) a(\alph) I(\Roman) i(\roman) 1(\arabic)
	%设定全局标号series=example	%引用全局变量resume=example
	%[topsep=-0.3em,parsep=-0.3em,itemsep=-0.3em,partopsep=-0.3em]
	%可使用leftmargin调整列表环境左边的空白长度 [leftmargin=0em]
	\item
	将上题中的区域$  D  $保形变换到上半平面。(6分)
	\item 
 利用上半平面 Laplace 方程第一边值问题的 Green 函数，写出 Laplace 方程第一边值问题在区域 $D$ 的 Green 函数，其中 $\bar{z}_{0}$ 是 $z_{0}$ 的共轭复数。 (2分)
	\[
	G\left(z ; z_{0}\right)=\frac{1}{2 \pi}\left(\ln \frac{1}{\left|z-z_{0}\right|}-\ln \frac{1}{|z-\bar{z}_{0}|}\right)
	\]
	
	
	
\end{enumerate}

\banswer{
	\begin{enumerate}
		\item
		写 $z=x+iy$。先令 $w=e^z=e^xe^{iy}$。由于
		$x<0$ 且 $0<y<\pi$，有
		\[
		|w|=e^x<1,\qquad \operatorname{Im}w>0,
		\]
		所以 $D$ 映为上半单位圆盘。其边界
		$y=0,\pi$ 映为直径 $(-1,1)$，而 $x=0$ 映为上半圆周。

		再作 Möbius 变换
		\[
		\eta=\frac{1+w}{1-w}.
		\]
		直径映到正实轴，上半圆周映到正虚轴，故上半单位圆盘
		映为第一象限。最后令 $\zeta=\eta^2$，第一象限便映为
		上半平面。因此复合映射为
		\[
		\zeta=\Phi(z)
		=\left(\frac{1+e^z}{1-e^z}\right)^2
		=\coth^2\frac z2.
		\]
		\item
		上半平面 $\mathbb H$ 的 Green 函数为
		\[
		G_{\mathbb H}(\zeta;\zeta_0)
		=\frac1{2\pi}
		\ln\left|
		\frac{\zeta-\overline{\zeta_0}}
		{\zeta-\zeta_0}
		\right|.
		\]
		这里它在实轴上为零，并且在 $\zeta=\zeta_0$ 附近具有
		$-(2\pi)^{-1}\ln|\zeta-\zeta_0|$ 型奇性。二维 Laplace
		方程的 Green 函数在保形映射下保持不变，所以取
		$\zeta=\Phi(z)$、$\zeta_0=\Phi(z_0)$，便有
		\[
		\boxed{\displaystyle
		G_D(z;z_0)=\frac1{2\pi}
		\ln\left|
		\frac{\Phi(z)-\overline{\Phi(z_0)}}
		{\Phi(z)-\Phi(z_0)}
		\right|}.
		\]
	\end{enumerate}
}

	
\end{enumerate}
